\documentclass[12pt]{article}
\usepackage{graphicx} % Required for inserting images
\usepackage[margin=0.5in]{geometry}
\usepackage{amsmath, amssymb}
\usepackage{mathrsfs}


\usepackage{soul}


% Dr. David Brown Commands
\newcommand{\ds}{\displaystyle}
\newcommand{\R}{\mathbb{R}}
\newcommand{\N}{\mathbb{N}}
\newcommand{\Q}{\mathbb{Q}}
\newcommand{\C}{\mathbb{C}}


% Commands to get script r
\usepackage{calligra}
\DeclareMathAlphabet{\mathcalligra}{T1}{calligra}{m}{n}
\DeclareFontShape{T1}{calligra}{m}{n}{<->s*[2.2]callig15}{}
\newcommand{\scripty}[1]{\ensuremath{\mathcalligra{#1}}}

% Unit commands
\newcommand{\degree}{\text{°}}
\newcommand{\unitSpeed}{\frac{\text{m}}{\text{s}}}
\newcommand{\unitAcceleration}{\frac{\text{m}}{\text{s}^2}}

% Constant commands
\newcommand{\avogadro}{6.022\times10^{23}}

\newcommand{\boltzmannConstK}{1.381\times10^{-23}\frac{\text{J}}{\text{K}}}
\newcommand{\boltzmannConstInEV}{8.617\times10^{-5}\frac{\text{eV}}{\text{K}}}
\newcommand{\planckConst}{6.626\times10^{-34}\text{J}\cdot\text{s}}

\newcommand{\atmP}{1.01325\times10^5\,\text{Pa}}

% Known Value Commands
\newcommand{\electronMass}{9.109\times10^{-31}\,\text{kg}}
\newcommand{\protonMass}{1.673\times10^{-27}\,\text{kg}}

% Commands to easily write partial derivatives
\newcommand{\partDeriv}[2]{\frac{\partial #1}{\partial #2}}
\newcommand{\doublePartDeriv}[2]{\frac{\partial^2 #1}{\partial^2 #2}}


\title{Problem 7.48}
\author{Gabriel Decker}


\begin{document}


\maketitle
\begin{flushleft}



In this problem, we will consider a neutrino-antineutrino gas.
We're given that neutrinos are fermions, and we will first approximate them to be massless.
We're also given the following chemical equation.
\begin{equation}
    \nu+\overline{\nu}\leftrightarrow2\gamma
\end{equation}
with $\nu$ representing a neutrino, $\overline{\nu}$ an antineutrino, and $\gamma$ a photon.
We also are given that $\mu_\nu=\mu_{\overline{\nu}}$.\\~\\


\textbf{Part a}\\
First, we must show that when the above chemical equation is in an equilibrium state, the chemical potentials $\mu_\nu=\mu_{\overline{\nu}}=0$.
Recall that $\mu_\gamma=0$, as shown early on in section 7.4.
Around page 210 (chapter 5), the book explains the conditions for when a system of different particle species is in chemical equilibrium.
From that, we know the following.
$$1\cdot\mu_\nu+1\cdot\mu_{\overline{\nu}}=2\cdot\mu_\gamma$$
$$\implies\mu_\nu+\mu_{\overline{\nu}}=2\mu_\gamma$$
$$\implies\mu_\nu+\mu_{\overline{\nu}}=0$$
$$\implies\mu_\nu=-\mu_{\overline{\nu}}$$
Since $\mu_\nu=\mu_{\overline{\nu}}$, this implies that they both equal zero.\\~\\~\\


\textbf{Part b}\\
We will now calculate the total energy per unit volume for this system.
We assume that the neutrinos are massless, relativistic fermions.
We're also given that there are three types of neutrinos that each have antineutrino species.
So, we have an extra degeneracy of $3\cdot2=6$.
This implies that we can use the same energy level equation that we did for the photon gas:
\begin{equation}
    \epsilon=\frac{hcn}{2L}
\end{equation}
with $n=\sqrt{n_x^2+n_y^2+n_z^2}$.
Recall the Fermi-Dirac distribution for Fermions.
\begin{equation}
    \overline{n}_\text{FD}=\frac{1}{e^{(\epsilon-\mu)/kT}+1}
\end{equation}\\~\\

Following a similar strategy as with the photon gas, the total energy is the following.
$$U=6\sum_\epsilon\epsilon\overline{n}_\text{FD}$$
$$\implies U=6\sum_\epsilon\frac{\epsilon}{e^{(\epsilon-\mu)/kT}+1}$$
Since $\mu=0$, we get the following.
$$\implies U=6\sum_\epsilon\frac{\epsilon}{e^{\epsilon/kT}+1}$$
Again, let's express the summation using $n$-space (like in problem 7.44 or the book's photon gas derivation).
$$\implies U=6\sum_{n_x}\sum_{n_y}\sum_{n_z}\frac{hcn}{2L}\frac{1}{e^{hcn/2LkT}+1}$$
Using the same $n$-space as in previous examples, we get the following integral in spherical coordinates.
$$\implies U=6\int_0^{\pi/2}\int_0^{\pi/2}\int_0^\infty\frac{hcn}{2L}\frac{1}{e^{hcn/2LkT}+1}n^2\sin(\theta)\,dn\,d\phi\,d\theta$$



\end{flushleft}

\end{document}
